\section*{अध्याय \ref{ch:b2:ism-star-formation} --- अंतरतारकीय माध्यम और तारों का निर्माण}

\begin{solution}{exo:b2:ism-star-formation:1}
$n = \num{2.7e-24}/\num{1.67e-24} = 1.6$ हाइड्रोजन परमाणु प्रति घन सेंटीमीटर। वायु में, आदर्श गैस के नियम से,
$n = P/k_{\mathrm B}T = 101325/(\num{1.381e-23}\times288) = \qty{2.55e25}{m^{-3}}$,
अर्थात् प्रति घन सेंटीमीटर \num{2.5e19} अणु: लगभग $1.6\times10^{19}$ गुना
अधिक।
\end{solution}

\begin{solution}{exo:b2:ism-star-formation:2}
$M_J \propto n^{-1/2}$: $n$ को 100 से भाग देने पर $M_J$ 10 से गुणा हो जाता है, इसलिए
$M_J = 54.6$ सौर द्रव्यमान, प्रति घन सेंटीमीटर $10^2$ कणों पर। मेघ में
$\num{24000}/54.6 \approx 440$ \omterm{def:b2:ism-star-formation:jeans}{जीन्स द्रव्यमान} हैं: वह समग्र रूप से स्थायी नहीं हो सकता, और उसे
विखंडित होना ही होगा।
\end{solution}

\begin{solution}{exo:b2:ism-star-formation:3}
$c_s = \sqrt{\num{1.381e-23}\times10/(2.33\times\num{1.661e-27})} =
\qty{189}{m.s^{-1}}$, \qty{10}{K} पर; $\sqrt{10}$ गुना अधिक, \qty{597}{m.s^{-1}},
\qty{100}{K} पर। \qty{0.1}{pc} पार करना: $\num{3.086e15}/189 = \qty{1.63e13}{s}$,
अर्थात् \qty{0.52}{Myr}।
\end{solution}

\begin{solution}{exo:b2:ism-star-formation:4}
$\rho = 2.33\times\num{1.661e-27}\times n$, जहाँ $n$ \unit{m^{-3}} में है। प्रति घन सेंटीमीटर $10^5$
कणों के लिए $\rho = \qty{3.87e-16}{kg.m^{-3}}$ और
$t_{\mathrm{ff}} = \sqrt{3\pi/(32G\rho)} = \qty{3.38e12}{s} = \qty{0.107}{Myr}$;
$10^2$ के लिए $\rho$ \num{1000} गुना छोटा है और $t_{\mathrm{ff}}$ $\sqrt{1000} = 31.6$
गुना लंबा, \qty{3.38}{Myr}। \omterm{def:b2:ism-star-formation:molecular-cloud}{सघन क्रोड} 32 गुना अधिक तेज़ी से ढहता है।
\end{solution}

\begin{solution}{exo:b2:ism-star-formation:5}
$k_{\mathrm B}T/(\mu m_u) = c_s^2$ लेने पर $(5k_{\mathrm B}T/G\mu m_u)^{3/2} =
5^{3/2}c_s^3G^{-3/2}$, जिससे सूत्र मिलता है। चूँकि $M_J \propto \rho^{-1/2}$ है,
$M_J$ आधा करने के लिए चार गुना घनत्व चाहिए। चूँकि $M_J \propto T^{3/2}$ है,
\qty{10}{K} से \qty{20}{K} तक गर्म होने पर वह $2^{3/2} = 2.83$ से गुणा हो जाता है।
\end{solution}

\begin{solution}{exo:b2:ism-star-formation:6}
$t_{\mathrm{KH}} = GM^2/(RL) = \num{6.674e-11}\times(\num{1.989e30})^2/(2\times
\num{6.96e8}\times2\times\num{3.83e26}) = \qty{2.48e14}{s}$, अर्थात्
\qty{7.8}{Myr}। यह अधिकांश चक्रिकाओं की आयु से लंबा है (लगभग \qty{3}{Myr} के बाद
आधी समाप्त हो जाती हैं): जब चक्रिका लुप्त होती है, तब तारा अब भी सिकुड़ रहा होता है।
\end{solution}

\begin{solution}{exo:b2:ism-star-formation:7}
\qty{10}{K} पर $c_s^3/G = 189^3/\num{6.674e-11} = \qty{1.01e17}{kg.s^{-1}}$, अर्थात्
$\num{1.01e17}\times\num{3.156e7}/\num{1.989e30} = \num{1.6e-6}$ सौर द्रव्यमान प्रति
वर्ष, और एक सौर द्रव्यमान में \qty{0.62}{Myr} लगते हैं। \qty{20}{K} पर दर
$2^{3/2} = 2.83$ गुना अधिक है, प्रति वर्ष \num{4.5e-6} सौर द्रव्यमान, और एक सौर
द्रव्यमान में \qty{0.22}{Myr} लगते हैं।
\end{solution}

\begin{solution}{exo:b2:ism-star-formation:8}
$r = \sqrt{L/(16\pi\sigma)}/T^2$, जहाँ $\sqrt{\num{3.83e26}/(16\pi\times\num{5.67e-8})}
= \qty{1.159e16}{m.K^2}$: \qty{1300}{K} पर $r = \qty{6.86e9}{m} = \qty{0.046}{au}$,
बुध की कक्षा के काफ़ी भीतर। \qty{5.2}{au} पर $T = 278/\sqrt{5.2} = \qty{122}{K}$।
\end{solution}

\begin{solution}{exo:b2:ism-star-formation:9}
$v = \sqrt{GM/r} = \sqrt{\num{6.674e-11}\times1.3\times\num{1.989e30}/\num{1.496e13}}
= \qty{3.40}{km.s^{-1}}$, और $P = 2\pi r/v = \qty{2.77e10}{s} = 877$ वर्ष
(तुल्यतः $\sqrt{100^3/1.3}$ वर्ष)। दस लाख वर्षों में गैस ने
$10^6/877 \approx \num{1140}$ चक्कर लगाए हैं: अपने घूर्णन की तुलना में चक्रिका पुरानी है,
इसीलिए उसके वलय चिकने वृत्त हैं।
\end{solution}

\begin{solution}{exo:b2:ism-star-formation:10}
$M = \int_0^R 4\pi r^2\rho_0R/r\,\mathrm dr = 2\pi\rho_0R^3$, और $r$ के भीतर
$m(r) = 2\pi\rho_0Rr^2$। $\mathrm dm = 4\pi\rho_0Rr\,\mathrm dr$ लेने पर,
\[
  U = -\int_0^R \frac{G\,m(r)}{r}\,\mathrm dm
    = -8\pi^2G\rho_0^2R^2\int_0^R r^2\,\mathrm dr = -\frac{8\pi^2}{3}G\rho_0^2R^5
    = -\frac23\,\frac{GM^2}{R},
\]
जहाँ $\rho_0 = M/(2\pi R^3)$ लिया गया है। यह एकसमान मान का $\tfrac{2/3}{3/5} = 1.11$ गुना है:
सांद्रित मेघ अधिक दृढ़ता से बँधा है और उसे थामे रखने के लिए \qty{11}{\percent} अधिक
तापीय ऊर्जा चाहिए; वह अधिक आसानी से ढहता है।
\end{solution}

\begin{solution}{exo:b2:ism-star-formation:11}
आवर्तकाल $2\pi/\Omega = \qty{6.28e14}{s}$, लगभग \qty{20}{Myr}।
$R = \qty{3.086e15}{m}$ के साथ $j = \Omega R^2 = \qty{9.52e16}{m^2.s^{-1}}$, और
$r_c = j^2/(GM) = \num{9.07e33}/\num{1.327e20} = \qty{6.83e13}{m} = \qty{457}{au}$।
दृढ़ घूर्णन में $j = \Omega r^2$ अक्ष से दूरी के वर्ग के अनुसार बढ़ता है:
क्रोड के भीतर की गैस में कोणीय संवेग कम होता है, और वह तारे के अधिक निकट उतरती है।
\end{solution}

\begin{solution}{exo:b2:ism-star-formation:12}
$\bar t = 2.46$, $\bar f = 0.588$; $\sum(t - \bar t)^2 = 12.612$,
$\sum(t - \bar t)(f - \bar f) = -2.018$, इसलिए $b = 2.018/12.612 = 0.160$ प्रति मिलियन
वर्ष और $a = 0.588 + 0.160\times2.46 = 0.982$। रेखा $0.5$ तक
$(0.982 - 0.5)/0.160 = \qty{3.0}{Myr}$ पर और शून्य तक $0.982/0.160 = \qty{6.1}{Myr}$ पर पहुँचती है, जैसा
चित्र में है। कोई अंश शून्य से नीचे नहीं जा सकता: पुराने तारागुच्छ, जैसे \qty{30}{Myr} का NGC 1960,
शून्य के निकट एक छोटा अंश बनाए रखते हैं जिसका वर्णन सरल रेखा नहीं कर सकती।
\end{solution}

\begin{solution}{pb:b2:ism-star-formation:1}
\textbf{1.} $R = 0.08\times\num{3.086e16} = \qty{2.469e15}{m}$।

\textbf{2.} $N = \qty{1.5e26}{m^{-2}}$, इसलिए $2.8\times\num{1.661e-27}\times\num{1.5e26} =
\qty{0.698}{kg.m^{-2}}$।

\textbf{3.} $M = 0.698\times\pi R^2 = 0.698\times\num{1.915e31} = \qty{1.34e31}{kg}$,
अर्थात् $6.7$ सौर द्रव्यमान।

\textbf{4.} $V = \tfrac43\pi R^3 = \qty{6.30e46}{m^3}$ और
$\rho = M/V = \qty{2.12e-16}{kg.m^{-3}}$।

\textbf{5.} $n = \rho/(2.33m_u) = \qty{5.48e10}{m^{-3}}$, अर्थात् प्रति घन सेंटीमीटर \num{5.5e4}
कण।

\textbf{6.} हाँ: उसका घनत्व $10^4$ और $10^5$ के बीच है, उसकी त्रिज्या 0.03 और
\qty{0.2}{pc} के बीच; उसका द्रव्यमान, 6.7 सौर द्रव्यमान, क्रोडों के परास के ऊपरी सिरे पर है।

\textbf{7.} $M_J = 5.46\times(10^4/\num{5.48e4})^{1/2} = 2.33$ सौर द्रव्यमान।

\textbf{8.} $6.7/2.33 = 2.9$ \omterm{def:b2:ism-star-formation:jeans}{जीन्स द्रव्यमान}: क्रोड अस्थायी है और ढहता है।

\textbf{9.} $R_J = 0.1316\times(10^4/\num{5.48e4})^{1/2} = \qty{0.056}{pc}$, जो
त्रिज्या, \qty{0.08}{pc}, से छोटी है: वही निष्कर्ष।

\textbf{10.} $2^{3/2}\times2.33 = 6.6$ सौर द्रव्यमान: लगभग क्रोड का द्रव्यमान।
\qty{20}{K} पर क्रोड सीमांत होता; उसका ठंडा होना ही उसे ढहाता है।

\textbf{11.} $t_{\mathrm{ff}} = \sqrt{3\pi/(32\times\num{6.674e-11}\times\num{2.12e-16})}
= \qty{4.56e12}{s}$, लगभग \num{145000} वर्ष।

\textbf{12.} $R/c_s = \num{2.469e15}/189 = \qty{1.31e13}{s}$, \num{414000} वर्ष: मुक्त पतन से
तीन गुना लंबा। दाब की तरंग क्रोड के ढहने से पहले उसे पार नहीं कर सकती;
दाब इतनी तेज़ी से प्रतिक्रिया नहीं कर सकता कि निपात को रोक दे।

\textbf{13.} $c_s = \qty{189}{m.s^{-1}}$ (\cref{exo:b2:ism-star-formation:3})।

\textbf{14.} $c_s^3/G = \qty{1.01e17}{kg.s^{-1}}$, प्रति वर्ष \num{1.6e-6} सौर द्रव्यमान।

\textbf{15.} $1/\num{1.6e-6} = \num{624000}$ वर्ष।

\textbf{16.} 6.7 का एक-तिहाई $2.2$ है: एक सौर द्रव्यमान के दो तारे।

\textbf{17.} $t_{\mathrm{KH}} = GM_\odot^2/(3R_\odot\times4L_\odot) = \qty{8.25e13}{s}$,
\qty{2.6}{Myr}।

\textbf{18.} लंबा, अभिवृद्धि-काल का चार गुना: जब अभिवृद्धि रुकती है, तब तारा
अब भी सिकुड़ रहा होता है और बिना संलयन के चमकता है; वह एक \omterm{def:b2:ism-star-formation:t-tauri}{टी टॉरी तारा} है, जो अभी मुख्य
अनुक्रम पर नहीं पहुँचा।

\textbf{19.} $\Omega = \qty{0.4e3}{m.s^{-1}}/\qty{3.086e16}{m} = \qty{1.30e-14}{rad.s^{-1}}$।

\textbf{20.} $j = 1.30\times10^{-14}\times(\num{2.469e15})^2 = \qty{7.90e16}{m^2.s^{-1}}$।

\textbf{21.} $r_c = j^2/(GM_\odot) = \num{6.24e33}/\num{1.327e20} = \qty{4.70e13}{m}$,
अर्थात् \qty{314}{au}: एचएल टॉरी की चक्रिका से बड़ी, जैसी कि क्रोड के किनारे की गैस से
अपेक्षा है।

\textbf{22.} $T = (L/16\pi\sigma r^2)^{1/4} = \qty{278}{K}$, \qty{1}{au} पर, और
$278/\sqrt{30} = \qty{51}{K}$, \qty{30}{au} पर।

\textbf{23.} $r = \num{1.159e16}/180^2 = \qty{3.58e11}{m}$, अर्थात् \qty{2.4}{au}, जैसा पाठ में है।

\textbf{24.} एक सौर द्रव्यमान के तारे को जुटने में
\textbf{\qty[text-series-to-math=true]{624}{} हज़ार वर्ष} लगते हैं, जो उसके क्रोड के
\omterm{def:b2:ism-star-formation:free-fall}{मुक्त पतन काल} (\num{145000} वर्ष) का लगभग चार गुना है: बाहरी गैस
\omterm{def:b2:ism-star-formation:protostar}{आद्यतारे} तक केवल धीरे-धीरे, भीतर से बाहर की ओर, पहुँचती है।
\end{solution}
